\section{Inverse mapping from electromagnetic fields to brane initial data}
\label{app:inverse-em-mapping}

In the main text, the electromagnetic sector is constructed in a
\emph{forward} direction: coarse-grained amplitude deformations of the
brane substrate source an effective scalar potential
$\Phi(\mathbf x) = \kappa\,\bar X^4(\mathbf x)$
(Eq.~\eqref{eq:phi-from-xi4}), which obeys a Poisson equation in the
quasi-static regime and is then promoted to the time component of a
four-potential $A_\mu$ in the Maxwell-type wave
equation~\eqref{eq:maxwell-wave}. In that construction the brane $\to$
EM map is regarded as primary.

For the purposes of initializing numerical experiments we require a
\emph{restricted inverse mapping} in the opposite direction: a local
dictionary that takes a classical electromagnetic field configuration
$(\mathbf E,\mathbf B)$ on the simulation patch and assigns compatible
microscopic initial data $(\mathbf X,\partial_t\mathbf X)$ for the
brane embedding. This appendix spells out such a mapping for the class
of approximately null, wave-like configurations used in the simulations.
It is \emph{not} a general inversion of Maxwell's equations, but a
physically motivated closure that matches local energy density and
energy flow between the EM and brane descriptions.

\subsection{Electromagnetic field configuration and effective-medium quantities}

We consider electromagnetic field configurations on a spatial patch,
which in the simplest case may be oriented along the $x$-axis with
transverse coordinates $(y,z)$ and spatial position
$\mathbf x = (x,y,z)$. The electromagnetic fields
\begin{equation}
  \mathbf E(\mathbf x,t), \qquad \mathbf B(\mathbf x,t),
\end{equation}
are described in a homogeneous effective medium characterized by
parameters $\varepsilon_{\mathrm{eff}}$ and $\mu_{\mathrm{eff}}$.
Consistency with the emergent Maxwell wave equation
(Eq.~\eqref{eq:maxwell-wave}) and the static Poisson limit fixes the
usual relation
\begin{equation}
  c^2
  \;=\;
  \frac{1}{\varepsilon_{\mathrm{eff}}\mu_{\mathrm{eff}}}\,,
  \label{eq:app-c2}
\end{equation}
which identifies $c$ with the wave speed of small disturbances in the
medium, in agreement with Eq.~\eqref{eq:c-squared-relation}.

The local energy density and Poynting vector of the effective
electromagnetic field are
\begin{equation}
  u_{\mathrm{EM}}(\mathbf x,t)
  = \frac{1}{2}\Bigl(
      \varepsilon_{\mathrm{eff}}\,|\mathbf E(\mathbf x,t)|^2
      + \frac{|\mathbf B(\mathbf x,t)|^2}{\mu_{\mathrm{eff}}}
    \Bigr),
  \qquad
  \mathbf S(\mathbf x,t)
  = \frac{1}{\mu_{\mathrm{eff}}}\,\mathbf E(\mathbf x,t)\times\mathbf B(\mathbf x,t).
  \label{eq:app-u-S}
\end{equation}
For a monochromatic null wave with $|\mathbf E| = c|\mathbf B|$ one has
$u_{\mathrm{EM}} = \varepsilon_{\mathrm{eff}}|\mathbf E|^2
                 = |\mathbf B|^2/\mu_{\mathrm{eff}}$
and
\begin{equation}
  \mathbf S
  = u_{\mathrm{EM}}\,\mathbf v_{\mathrm{energy}},
  \qquad
  |\mathbf v_{\mathrm{energy}}| \approx c,
\end{equation}
where $\mathbf v_{\mathrm{energy}}$ is the usual energy-transport
velocity of the guided mode, aligned with the waveguide axis in the
straight case.

\subsection{Amplitude mode and local energy matching}
\label{app:inverse-amplitude}

On the microscopic side, the brane dynamics are governed by the
Lagrangian density~\eqref{eq:brane-lagrangian},
\begin{equation}
  \mathcal{L}
  = \frac{\rho_m}{2}|\partial_t\mathbf X|^2
    - \frac{T}{2}\mathrm{tr}(E)
    - \mu |E|^2
    - \frac{\kappa}{2}|b|^2,
\end{equation}
where $\rho_m$ is the mass density and $T$ the effective tension of the
medium. In the static gauge used throughout the paper we write
\begin{equation}
  X^a(x,t) = x^a, \qquad X^4(x,t) = \xi_\perp(x,t),
\end{equation}
so that $\xi_\perp$ denotes the normal displacement (amplitude) of the
brane in the fourth embedding direction.

For small, predominantly longitudinal deformations along the waveguide
we can approximate the dynamics of $\xi_\perp$ by an effective
one-dimensional wave equation along the guide coordinate $s$,
\begin{equation}
  \rho_m\,\partial_{tt}\xi_\perp(s,t)
  \;=\;
  T\,\partial_{ss}\xi_\perp(s,t),
  \qquad
  c^2
  \;=\;
  \frac{T}{\rho_m},
  \label{eq:app-1d-wave}
\end{equation}
with local energy density
\begin{equation}
  u_{\mathrm{brane}}(s,t)
  = \frac{1}{2}\rho_m\bigl(\partial_t\xi_\perp\bigr)^2
    + \frac{1}{2}T\bigl(\partial_s\xi_\perp\bigr)^2.
  \label{eq:app-u-brane}
\end{equation}

In the linear regime the normal displacement $\xi_\perp$ obeys the
massless one-dimensional wave equation~\eqref{eq:app-1d-wave}. Any
solution can locally be decomposed into left- and right-moving wave
packets $\xi_\perp(s,t) = f(s-ct) + g(s+ct)$. For such traveling waves
one has the local relation
\begin{equation}
  \partial_t\xi_\perp = \mp c\,\partial_s\xi_\perp,
\end{equation}
so that the energy density~\eqref{eq:app-u-brane} simplifies to
\begin{equation}
  u_{\mathrm{brane}}(s,t)
  = \rho_m c^2\bigl(\partial_s\xi_\perp\bigr)^2
  = \frac{\rho_m}{c^2}\bigl(\partial_t\xi_\perp\bigr)^2.
\end{equation}
This shows that in the traveling-wave regime the brane energy density
is a quadratic form in the spatial or temporal derivative of
$\xi_\perp$, without reference to a particle picture or a specific
oscillation phase.

For the purposes of the inverse mapping we compress the above relation
into a single effective calibration constant $K_{\mathrm{eff}}$ and
write, schematically,
\begin{equation}
  u_{\mathrm{brane}}(s,t)
  \;\approx\;
  \frac{1}{2} K_{\mathrm{eff}}\,\Xi^2(s,t),
\end{equation}
where $\Xi$ is some chosen local combination of $\partial_t\xi_\perp$
and $\partial_s\xi_\perp$ (for instance the derivative along the
propagation direction). The constant $K_{\mathrm{eff}}$ is determined
once and for all from the microscopic brane parameters $(\rho_m,T,\dots)$
and the Compton-scale calibration discussed in the main text. It is not
an additional dynamical parameter and in particular it does not encode
any ``photon frequency''.

We now demand local energy matching between the EM and brane
descriptions on the patch,
\begin{equation}
  u_{\mathrm{brane}}(\mathbf x,0)
  \;\approx\;
  u_{\mathrm{EM}}(\mathbf x,0),
\end{equation}
where
\begin{equation}
  u_{\mathrm{EM}}(\mathbf x,0)
  = \frac{1}{2}\Bigl(
      \varepsilon_{\mathrm{eff}}\,|\mathbf E(\mathbf x,0)|^2
      + \frac{|\mathbf B(\mathbf x,0)|^2}{\mu_{\mathrm{eff}}}
    \Bigr)
  \label{eq:app-u-EM}
\end{equation}
is the usual electromagnetic energy density.

Using the effective quadratic form for the brane energy we write
\begin{equation}
  u_{\mathrm{brane}}(\mathbf x,0)
  = \frac{1}{2} K_{\mathrm{eff}}\,\Xi^2(\mathbf x,0),
\end{equation}
and choose a simple gauge in which all the energy is initially stored
in the ``potential'' part of the normal deformation. Concretely we set
\begin{equation}
  \partial_t\xi_\perp(\mathbf x,0) = 0,
\end{equation}
and define
\begin{equation}
  \xi_\perp(\mathbf x,0)
  = A(\mathbf x)
  = \sqrt{\frac{2\,u_{\mathrm{EM}}(\mathbf x,0)}{K_{\mathrm{eff}}}}.
  \label{eq:app-xi-initial}
\end{equation}

The choice $\partial_t\xi_\perp(\mathbf x,0)=0$ fixes a convenient
phase convention for the microscopic brane field. It does not affect
the coarse-grained electromagnetic configuration, which depends only on
the combination of derivatives entering $u_{\mathrm{brane}}$ and the
emergent Maxwell sector. In particular, no additional ``oscillator
phase'' field is introduced in the inverse mapping; the only inputs
from the EM side are the classical fields $\mathbf E$ and $\mathbf B$
via $u_{\mathrm{EM}}$.

In numerical experiments we typically identify $\xi_\perp$ with the
fourth embedding component $X^4$ in the small-slope regime,
$\xi_\perp\approx X^4$, and restrict $A$ to remain within the linear
elastic range, so that higher-order geometric nonlinearities are
controlled by the subsequent evolution rather than the initial
conditions.

\subsection{Energy--momentum matching and lateral velocities}
\label{app:poynting-lateral}

In the ontology of this model, magnetic effects are not independent
degrees of freedom but arise as \emph{relativistic manifestations} of
electric fields in motion. At the level of the effective electromagnetic
description this is encoded in the field tensor $F_{\mu\nu}$ and the
Lorentz transformation properties of $(\mathbf E,\mathbf B)$. At the
microscopic level we instead interpret the energy flow associated with
$\mathbf B$ as a statement about the \emph{lateral motion} of brane
material.

At the level of the energy--momentum tensor, the electromagnetic
momentum density is $T^{0i}_{\mathrm{EM}} = S^i/c^2$. In our brane
model, lateral motion of the embedding carries momentum. For the class
of null, wave-like configurations considered here we therefore identify
the lateral velocity of the brane with the local energy-flow velocity
of the electromagnetic field,
\begin{equation}
  \partial_t\mathbf X_\parallel(\mathbf x,0)
  = \mathbf v_{\mathrm{energy}}(\mathbf x,0)
  = \frac{\mathbf S(\mathbf x,0)}{u_{\mathrm{EM}}(\mathbf x,0)}.
  \label{eq:app-Xparallel-initial}
\end{equation}
This identification is formulated at the level of fields and their
energy--momentum tensors and does not invoke a particle interpretation.
The normal velocity $\partial_t X^4(\mathbf x,0)$ is already fixed by
Eq.~\eqref{eq:app-xi-initial} to be zero in our gauge choice. In regions
where $u_{\mathrm{EM}}$ is negligible, the ratio in
Eq.~\eqref{eq:app-Xparallel-initial} is suppressed and the brane remains
approximately at rest in the tangential directions, as expected for
vacuum regions outside the wave packet.

In a Lorentz-covariant language one may view this prescription as
selecting, at each point on the patch, a local frame in which the
field is predominantly electric and relating the observed magnetic
component to the lateral motion of the underlying amplitude deformation.
Here we bypass the explicit frame transformation and use the Poynting
vector directly as a proxy for the microscopic energy flow.

\subsection{Extension to the toroidal double loop and limitations}

The construction above is local in $\mathbf x$ and does not rely on the
global geometry of the underlying configuration. It therefore carries
over to the toroidal electron picture of
Appendix~\ref{app:tubular-electron} without modification: one replaces
the coordinate $s$ by the longitudinal coordinate $z$ along the
centerline $\mathbf C(z)$ of the double loop, uses the tubular embedding
$\mathbf r(z,x,y)$ of Eq.~\eqref{eq:tubular-embedding}, and applies
Eqs.~\eqref{eq:app-xi-initial}--\eqref{eq:app-Xparallel-initial}
pointwise along the strip. The only additional requirement is that the
local curvature of the strip remain small compared to characteristic
length scales of field variation, so that the local plane-wave
approximation leading to Eq.~\eqref{eq:app-1d-wave} remains valid.

It is important to emphasize the limited scope of this inverse mapping:

\begin{itemize}
  \item It applies to approximately null, wave-like field configurations
    in an effective medium characterized by
    $\varepsilon_{\mathrm{eff}}$ and $\mu_{\mathrm{eff}}$ obeying
    Eq.~\eqref{eq:app-c2}. It is not intended as a general inversion
    procedure for arbitrary electromagnetic field configurations or
    superpositions.

  \item On the electromagnetic side, the only inputs to the inverse
    mapping are the fields $\mathbf E(\mathbf x,0)$ and
    $\mathbf B(\mathbf x,0)$ via the derived quantities
    $u_{\mathrm{EM}}$ and $\mathbf S$ in Eq.~\eqref{eq:app-u-S}. No
    additional parameters such as a prescribed ``photon frequency'' or
    guided-mode index are required.

  \item On the brane side, the mapping depends only on the microscopic
    material parameters $(\rho_m,T,\dots)$ through the effective constant
    $K_{\mathrm{eff}}$. These are fixed once by calibration and are not
    interpreted as external control parameters of the electromagnetic
    configuration.

  \item The mapping is used as a \emph{closure prescription} for
    constructing initial brane configurations in numerical experiments.
    The emergent electromagnetic fields at later times are still
    obtained by the \emph{forward} brane $\to$ EM dictionary of the main
    text, via $\bar X^4$, the Poisson relation and the Maxwell-type wave
    equation~\eqref{eq:maxwell-wave}.
\end{itemize}

Within these limitations, the inverse mapping provides a concrete and
internally consistent way to embed classical electromagnetic field
configurations (such as localized null wave packets) into the
microscopic brane dynamics, using only field-theoretic quantities and
without invoking a particle ontology at the mapping level.